\chapter{Il modello Graphagate: Temporal Graph Network per la Zero Trust Architecture}

\section{Introduzione}

Nel contesto di una Zero Trust Architecture (ZTA), il processo di autorizzazione continua richiede meccanismi dinamici per rilevare deviazioni dal comportamento abituale e segnali di compromissione. Il modello integrato con l'Orchestrator OPA è basato su una \textit{Temporal Graph Network} (TGN), progettata specificamente per il \textit{self-supervised anomaly detection} su stream di traffico continuo.

Ogni accesso (una richiesta IP/device verso una risorsa) rappresenta un arco temporale all'interno di un grafo in evoluzione. Il modello mantiene uno stato ricorrente per ciascuna entità (utente o risorsa) e analizza i vicini temporali recenti per assegnare uno score di anomalia in tempo reale, consentendo di bloccare le richieste ostili ed esfiltranti.

\section{Il modello base: Temporal Graph Network (v3)}
\label{sec:v3}

\subsection{Architettura del sistema}

L'architettura del modello combina la memoria ricorrente tipica dei TGN (Rossi et al., 2020) con un gestore del vicinato ottimizzato (bounded in RAM) e uno \textit{scorer} a due teste che permette di analizzare le anomalie su due fronti: contestuale/policy e strutturale.

\subsubsection{Componenti principali}

\begin{itemize}
    \item \textit{TGNMemory}: struttura basata su GRU che aggiorna lo stato storico di ciascun nodo con i messaggi degli eventi in ingresso. La dimensione è impostata a 256 per gestire le impronte comportamentali.
    \item \textit{Hashed Identity}: identità dei nodi mappate tramite funzioni di hash (\texttt{hash(URI) \% hash\_buckets}) per garantire totale \textit{induttività} e scalabilità in presenza di entità mai viste prima.
    \item \textit{MessageNeighborLoader}: buffer in RAM che memorizza gli ultimi vicini temporali di ciascun nodo (es. \texttt{neighbor\_size=30}). Lavora con lookup in tempo $O(1)$ su matrici preallocate ed elimina la necessità di lenti database a grafo.
    \item \textit{GraphAttentionEmbedding}: rete multi-hop con \textit{TransformerConv} per calcolare l'embedding attendendo sui vicini storici, concatenando l'identità alla storia temporale.
    \item \textit{Static Node Features}: vettore a 16 dimensioni associato a ciascun nodo, che codifica attributi ZTA essenziali come \textit{device tier}, \textit{clearance} e \textit{trust score}, oltre a metriche estese introdotte nello schema v3 quali il \textit{resource risk} (per innalzare la sensibilità su asset critici) e il \textit{source internal/external flag} (per distinguere reti RFC1918 da IP pubblici; ablatabile e \emph{disattivato di default} nella configurazione corrente, poiché il roaming esterno benigno può mascherare il furto di credenziali --- la coppia external$\rightarrow$device che esporrebbe l'attaccante).
    \item \textit{History Features} (esplicite, non circolari): per ogni arco valutato si calcolano tre contatori causali e \textit{benign-gated} (aggiornati solo su traffico approvato, anti-poisoning): $[\log(1{+}\text{pair\_count}),\ \log(1{+}\text{src\_count}),\ \text{pair\_count}/(\text{src\_count}{+}1)]$. A questi si aggiunge una seconda tripletta identica per le coppie \textit{aux\_src}$\rightarrow$\textit{dst} (es. per validare la \textit{device habituality}), portando la dimensione a 6. Misurano quanto è abituale quell'accesso: una coppia mai vista da una sorgente attiva è la \textit{novelty cue} del movimento laterale, che, combinata con la memoria temporale, porta il maggior contributo (ablation multi-seed a 3 seed: $+0.163$ AUC, il segnale dominante del laterale).
    \item \textit{Dual Head Scorer}: lo score di un arco è la \emph{somma} di due logit complementari, $\text{logit} = \text{feat\_logit} + \text{struct\_logit}$.
    \begin{itemize}
        \item \textit{Feature head} (\texttt{LinkPredictor}, MLP a 3 strati): consuma gli embedding dei due endpoint, il messaggio dell'evento, gli attributi statici con hashed-identity, le codifiche di recency ($\Delta t$ coppia, $\Delta t$ sorgente) e le \textit{history features}. Cattura anomalie \textit{contestuali} e di \textit{policy} e veicola la novelty cue del laterale.
        \item \textit{Structural head}: similarità coseno (scalata da una temperatura apprendibile) delle proiezioni non-lineari degli embedding. \emph{Nota}: l'ablation a 3 seed la mostra \emph{marginale} sul laterale ($0.882$ con vs $0.875$ senza, $+0.007$ AUC); su questo stream il segnale è già portato dalle \textit{history features} e dalla memoria temporale. La testa è comunque \emph{mantenuta} come segnale complementare a basso costo: l'ottimizzatore le assegna una temperatura non banale (\texttt{struct\_scale}$\approx 4.7$ nel checkpoint addestrato), dunque è \emph{attiva ma ridondante}, non inutilizzata. Poiché nel task sintetico il movimento laterale coincide quasi per costruzione con la \textit{pair-novelty} - catturata direttamente dalle history features - l'ablation ne sottostima probabilmente il contributo: il suo valore su topologie reali (dove il laterale ha struttura di grafo oltre alla semplice novità della coppia) e in regime di \textit{cold-start} o di \textit{poisoning} dei contatori, dove il coseno sugli embedding è un segnale indipendente dai contatori espliciti, è plausibile ma \emph{non ancora verificato}.
    \end{itemize}
    \item \textit{Kill-chain Precursor} (prior di serving, non addestrato): il movimento laterale segue tipicamente un alert di recon (Snort / anomalia rilevata) sulla \emph{stessa} entità, ma il gate predict-then-update scarta quel precursore dalla memoria TGN. Si mantiene quindi uno stato per-entità \texttt{recent\_alert} e si applica un fattore \emph{moltiplicativo} allo score, $1 + \text{max\_boost}\cdot 0.5^{\,\Delta t / \text{half\_life}}$ (default $\text{max\_boost}{=}2$, $\text{half\_life}{=}600$\,s). Moltiplicativo per costruzione: alza uno score già sospetto sopra soglia senza trasformare in falsi positivi gli score benigni $\approx 0$. Il suo contributo è però \emph{marginale} (ablation multi-seed a 3 seed: $+0.013$ AUC), comparabile a quello della testa strutturale. Non costituisce un input addestrato (il training solo-benigno lo renderebbe una feature morta).
\end{itemize}

\begin{figure}
\centering
\begin{tikzpicture}[
    font=\sffamily\scriptsize,
    inp/.style={draw, rounded corners, align=center, fill=green!10, inner sep=5pt, minimum height=0.6cm},
    op/.style={draw, rounded corners, align=center, fill=blue!5, inner sep=5pt, minimum height=0.55cm},
    hop/.style={draw, rounded corners, align=center, fill=blue!12, inner sep=4pt, text width=5.6cm, minimum height=0.5cm},
    lay/.style={draw, rounded corners, align=center, fill=orange!16, inner sep=4pt, text width=4.7cm, minimum height=0.45cm},
    note/.style={align=center, font=\sffamily\scriptsize\itshape},
    gaebox/.style={draw, rounded corners, fill=blue!4, inner sep=12pt},
    headbox/.style={draw, rounded corners, fill=orange!6, inner sep=8pt},
    scorerbox/.style={draw, densely dashed, rounded corners, fill=orange!2, inner sep=22pt},
    output/.style={draw, rounded corners, align=center, fill=red!8, inner sep=5pt, text width=11cm},
    lbl/.style={font=\sffamily\scriptsize\bfseries},
    fl/.style={->, >=Stealth, thick, rounded corners=3pt},
    res/.style={->, >=Stealth, thick, densely dashed, rounded corners=3pt}
]

% ===== RIGA INPUT =====
\node[inp] (NF) {nf $[16]$\\\textit{tier / risk / internal}};
\node[inp, left=1.1cm of NF] (ZMEM) {z\_mem $[256]$\\\textit{GRU TGNMemory}};
\node[inp, right=1.1cm of NF] (HE) {he $[16]$\\\textit{hash\_emb (100k buckets)}};

\node[op, below=0.9cm of NF] (CAT) {concat\quad x $[256{+}16{+}16 = 288]$};
\draw[fl] (ZMEM.south) |- (CAT.west);
\draw[fl] (NF.south) -- (CAT.north);
\draw[fl] (HE.south) |- (CAT.east);

% ===== GRAPH ATTENTION EMBEDDING: i 3 hop come sottocomponenti =====
\node[hop, below=1.5cm of CAT] (HOP1)
  {\textbf{hop 1}: TransformerConv$(288{\to}256)$\\$\to$ LayerNorm $\to$ ReLU};
\node[hop, below=0.5cm of HOP1] (HOP2)
  {\textbf{hop 2}: TransformerConv$(256{\to}256)$ $+$ \textit{residuo}\\$\to$ LayerNorm $\to$ ReLU};
\node[hop, below=0.5cm of HOP2] (HOP3)
  {\textbf{hop 3}: TransformerConv$(256{\to}256)$ $+$ \textit{residuo}\\$\to$ LayerNorm \textit{(no ReLU)} $\;\to\;$ \textbf{z}$[256]$};

\node[note, text width=4.0cm, left=1.0cm of HOP2] (EDGE)
  {\textbf{edge\_attr}\\$[\,$time\_enc$(\Delta t_{hist})\,\|\,$hist\_msg$\,]$\\$=[32{+}7]=39$};

\draw[fl] (CAT.south) -- (HOP1.north);
\draw[fl] (HOP1) -- (HOP2);
\draw[fl] (HOP2) -- (HOP3);
\draw[fl, densely dashed] (EDGE.east) -- (HOP2.west);
% connessioni residue (lato destro)
\draw[res] (HOP1.east) -- ++(0.85,0) |- (HOP2.east);
\draw[res] (HOP2.east) -- ++(0.85,0) |- (HOP3.east);
\coordinate (GAER) at ([xshift=1.0cm]HOP3.east);

% ===== DUAL HEAD SCORER: due teste, ciascuna coi propri layer =====
% --- Feature head ---
\node[lay, anchor=north] (FIN) at ([xshift=-4.0cm, yshift=-2.0cm]HOP3.south)
  {input $[653]=$ z\_src$\|$z\_dst$\|$msg$[7]$\\$\|\,$fwh\_src$[32]\,\|\,$fwh\_dst$[32]$\\$\|\,\Delta t[32]\,\|\,\Delta t_{src}[32]\,\|\,$hist$[6]$};
\node[lay, below=0.4cm of FIN] (F1) {Linear$(653{\to}256)$ + ReLU};
\node[lay, below=0.35cm of F1] (F2) {Linear$(256{\to}256)$ + ReLU};
\node[lay, below=0.35cm of F2] (F3) {Linear$(256{\to}1)\;\to\;$ feat\_logit};
\draw[fl] (FIN) -- (F1);
\draw[fl] (F1) -- (F2);
\draw[fl] (F2) -- (F3);

% --- Structural head ---
\node[lay, anchor=north] (S1) at ([xshift=4.0cm, yshift=-2.0cm]HOP3.south)
  {struct\_proj: Linear$(256{\to}512)$\\+ ReLU + Dropout$(0.1)$};
\node[lay, below=0.4cm of S1] (S2) {Linear$(512{\to}256)$\\\textit{(su z\_src e z\_dst)}};
\node[lay, below=0.35cm of S2] (S3) {struct\_scale$\cdot\cos(\hat z_{src},\hat z_{dst})$\\$\to$ struct\_logit \textit{(marginale)}};
\draw[fl] (S1) -- (S2);
\draw[fl] (S2) -- (S3);

% z -> le due teste
\draw[fl] (HOP3.south) -- ++(0,-0.7) -| (FIN.north);
\draw[fl] (HOP3.south) -- ++(0,-0.7) -| (S1.north);

% ===== MERGE + SIGMOID + PRECURSORE =====
\coordinate (BOT) at ([yshift=-1.1cm]F3.south);
\node[op, below=0.1cm] (SUM) at (HOP3 |- BOT) {logit $=$ feat\_logit $+$ struct\_logit};
\draw[fl] (F3.south) |- (SUM.west);
\draw[fl] (S3.south) |- (SUM.east);

\node[op, below=0.7cm of SUM] (SIG) {$P(\text{benign})=\sigma(\text{logit})$\quad$\rightarrow$\quad score $= 1-P(\text{benign})$};
\draw[fl] (SUM.south) -- (SIG.north);

\node[output, below=0.7cm of SIG] (PREC)
  {$\times$ \textbf{precursor\_boost} $= 1 + 2\cdot 0.5^{\,\Delta t/600}$ \textit{(prior di serving, non addestrato)}\\
   $\rightarrow$ score finale $=\min(1,\,\cdot)$, decisione instradata sul segnale (soglia clean / dirty)};
\draw[fl] (SIG.south) -- (PREC.north);

% ===== RIQUADRI DEI COMPONENTI D'ALTO LIVELLO (background) =====
\begin{scope}[on background layer]
  \node[gaebox, fit=(HOP1)(HOP2)(HOP3)(EDGE)(GAER),
        label={[lbl, anchor=south west]north west:GraphAttentionEmbedding - \texttt{neighbor\_size}=30}] {};
  % outer scorer prima (piano piu' arretrato), poi le due teste interne
  \node[scorerbox, fit=(FIN)(F1)(F2)(F3)(S1)(S2)(S3),
        label={[lbl, anchor=south west]north west:Dual Head Scorer}] (SCORER) {};
  \node[headbox, fit=(FIN)(F1)(F2)(F3),
        label={[lbl, anchor=south west]north west:Feature head --- LinkPredictor (MLP)}] {};
  \node[headbox, fit=(S1)(S2)(S3),
        label={[lbl, anchor=south west]north west:Structural head}] {};
\end{scope}

\end{tikzpicture}
\caption{Schema dei layer dove i moduli interni sono esplicitati per i due di alto livello (GAE ed il Dual Head Scorer): i tre hop dentro \textit{GraphAttentionEmbedding} e le due teste (\textit{Feature}/\textit{Structural}) dentro il \textit{Dual Head Scorer}. Tutti i \texttt{TransformerConv} usano \texttt{heads}=4, \texttt{dropout}=0.1 e concatenano \texttt{edge\_attr} (alimenta ogni hop). Lo score grezzo $1-P(\text{benign})$ è moltiplicato dal prior di kill-chain (con cap a 1) e confrontato con la soglia instradata sul segnale; a serving lo score dell'evento è il massimo sui suoi archi (accesso e binding).}
\label{fig:layers}
\end{figure}


\subsection{Fase di addestramento e calibrazione}

Il modello viene addestrato con un approccio \textit{self-supervised} unicamente su traffico benigno, senza mai essere esposto a etichette di anomalia (unsupervised anomaly detection). Lo stream temporale sintetico generato viene suddiviso cronologicamente in tre subset: addestramento (70\%), validazione (10\%) e test (20\%).

L'obiettivo di apprendimento per la rete è il \textit{link prediction}: il TGN è istruito per massimizzare lo score assegnato ad eventi realmente accaduti nella baseline benigna e minimizzare contemporaneamente lo score associato ad eventi artificialmente manipolati e generati a runtime (\textit{negative sampling}).

\subsubsection{Negative sampling e funzione di loss}

Il processo di ottimizzazione (guidato dall'algoritmo AdamW) somma tre termini self-supervised. La componente dominante è una loss di ranking InfoNCE, allineata all'Average Precision, affiancata da due ancoraggi BCE:

\begin{itemize}
    \item \textit{Structural Negatives} (InfoNCE, $K{=}5$ random target): per ogni positivo benigno la sorgente è accoppiata a $K$ destinazioni \emph{uniformemente casuali} sull'intero id-range delle destinazioni. L'obiettivo InfoNCE impone che, tra $\{$dst-vera, $K$ dst-casuali$\}$, la dst-vera ottenga il logit più benigno \emph{dato lo storico della sorgente}. La rete apprende così $P(\text{dst}\mid\text{src, storia})$: un accesso a bassa verosimiglianza (movimento laterale o violazione di policy) riceve uno score elevato.
    \item \textit{Positive BCE Anchor}: i logit dei positivi benigni vengono spinti verso ``benigno''. L'InfoNCE fissa solo l'ordine relativo; questo ancoraggio mantiene gli score su una scala assoluta, così che la soglia calibrata sull'FPR resti significativa.
    \item \textit{Contextual Negatives} (Gaussian Feature Noise): mantenendo intatta la destinazione, al messaggio dell'evento si somma rumore gaussiano ($\sigma{=}0.5$) e si addestra la \textit{Feature Head} a marcarlo come anomalo. Si tratta di un meccanismo distinto rispetto alle anomalie contestuali della valutazione (randomizzazione discreta dei segnali 0/1), così il modello non riceve la corruzione del test.
\end{itemize}

% \paragraph{De-circolarizzazione (nessun data leak nei negativi).} I negativi strutturali usano \emph{solo} l'id-range delle destinazioni, mai la matrice degli autorizzati né i set abituali del generatore. Una versione precedente impiegava un ``hard negative'' ristretto agli \textit{autorizzati-ma-non-abituali} (con peso $\times 10$): essendo \emph{la definizione esatta} del movimento laterale di test, produceva un recall circolare e gonfiato. È stato rimosso di proposito (vedi \texttt{\_sample\_structural\_negatives}); ogni miglioria deriva da architettura/segnale, mai da negativi che rispecchiano la regola di test. Inoltre il trust (\texttt{node\_feat[:,14]}) non viene più mutato dalle etichette in training, evitando di accoppiare la label in una feature di input.

\subsubsection{Calibrazione della soglia (cost-sensitive instradata sul segnale)}
\label{sec:threshold}

Il TGN non emette una classificazione netta ma uno score continuo confrontato con una soglia, calibrata via \textit{replay} sul 10\% di validazione held-out (lo split di test non è mai toccato in calibrazione). Poiché a serving la classe vera è ignota ma il \emph{segnale di edge} (JA3/Snort/sensori) è osservabile, si calibrano e si instradano due soglie:

\begin{itemize}
    \item \textit{Soglia signal-dirty} (eventi il cui segnale già scatta): il classico quantile al Target FPR dell'1\% sugli score benigni. Questi eventi sono già intercettati dalla cheap rule baseline, quindi si mantiene la soglia conservativa.
    \item \textit{Soglia signal-clean} (dove il laterale è indistinguibile dal benigno se non per il pattern temporale): soglia \emph{cost-sensitive} che minimizza $\text{cost\_ratio}\cdot\text{FN} + \text{FP}$ (default \texttt{cost\_ratio}=20: un laterale mancato costa molto più di un falso allarme ri-sfidabile dall'orchestrator). Un \emph{cap} esplicito sull'FPR del flusso clean (\texttt{clean\_fpr\_cap}=0.05) impedisce che la spinta sul recall faccia esplodere i falsi positivi --- necessario perché anche con lateral AUC$\sim$0.91 (3 seed) il trade-off recall/FPR sul flusso clean resta ripido (a recall laterale $\sim$79\% corrisponde un FPR benigno $\sim$13\%).
\end{itemize}

A inferenza la decisione è instradata per-evento: \texttt{signal\_dirty} $\Rightarrow$ soglia dirty, altrimenti soglia clean (cost-sensitive). Entrambe le soglie sono persistite ed esposte da \texttt{serve\_api} (anche su \texttt{/health}). OPA resta il decisore finale allow/deny.

\subsection{Tipologie di anomalie rilevate in produzione}

In fase di valutazione o in ambiente live, l'architettura rileva 3 macro-categorie di attacco, dimostrando vantaggi misurabili rispetto a baseline più semplici (rule-based, Isolation Forest, One-Class SVM e GNN non temporale; cfr. Sezione~\ref{sec:eval}):

\begin{enumerate}
    \item \textit{Policy Anomaly}: l'entità agisce al di fuori del proprio ruolo, clearance o device tier. Intercettata agevolmente dalla Feature head grazie alle feature statiche ZTA.
    \item \textit{Contextual Anomaly}: la connessione ha caratteristiche insolite o ha innescato sensori IPS. Rilevata dalla Feature head decostruendo le \textit{edge features}.
    \item \textit{Lateral Movement}: l'accesso è formalmente autorizzato a livello di regole OPA, ma \textit{non abituale}. Avviene tipicamente nel ``pivoting''. Essendo indistinguibile sintatticamente dal traffico benigno, la classificazione ricade quasi del tutto sull'architettura implementata.
\end{enumerate}

\subsection{Valutazione sperimentale e confronto con le baseline}
\label{sec:eval}

La valutazione è condotta sullo split di test (20\% finale dello stream sintetico) tramite replay per-evento attraverso l'esatto codice di serving (\texttt{num\_events}=200\,000, di cui il 20\% finale $=40\,000$ eventi di test), su \textbf{3 seed} \texttt{[42,7,123]} (media $\pm$ dev.std). Per quantificare il contributo \emph{effettivo} della componente temporale, il TGN è confrontato con tre baseline non supervisionate addestrate sullo \emph{stesso} traffico benigno, con identico split cronologico, identici segnali tabellari (contatori di storia + prior precursore), soglia calibrata all'1\% di FPR e protocollo di valutazione (codice in \texttt{tests/baselines/}); una quarta baseline \emph{supervisionata} (XGBoost) è riportata solo come upper-bound di riferimento:

\begin{itemize}
    \item \textit{Isolation Forest}: detector non relazionale sulle sole feature statiche per-evento (nessuna struttura di grafo) --- rappresenta il ``pavimento'' privo di informazione strutturale (lateral AUC $0.645\pm0.025$).
    \item \textit{One-Class SVM}: controparte kernel non relazionale dell'Isolation Forest (sole feature per-evento). Pur con AUC aggregata competitiva ($0.801\pm0.018$), sul \textit{lateral} resta a $0.599\pm0.002$, confermando che il segnale temporale non è catturabile dalle sole feature statiche.
    \item \textit{GNN non temporale}: ablation del TGN su grafo statico aggregato, con lo \emph{stesso} curriculum di negative sampling de-circolarizzato (InfoNCE su destinazioni casuali + rumore contestuale) e gli \emph{stessi} contatori tabellari di storia, ma priva di memoria ricorrente e vicinato temporale. Isola la sola componente temporale (lateral AUC $0.451\pm0.021$, $\approx$ caso).
    \item \textit{XGBoost} (\emph{supervisionato}, fuori categoria): addestrato con le etichette, raggiunge lateral AUC $0.929\pm0.001$. Non è un confronto equo col paradigma non supervisionato del TGN; è incluso solo per delimitare l'upper-bound raggiungibile avendo accesso alle label (assenti per costruzione nel nostro setting).
\end{itemize}

I risultati del confronto sono riportati nella Tabella~\ref{tab:baselines}.

\begin{table}[ht!]
\centering
\footnotesize
\renewcommand{\arraystretch}{1.4}
\setlength{\tabcolsep}{5pt}
\begin{tabular}{l c c c c}
\hline
\textbf{Metrica (test set)} & \textbf{TGN} & \textbf{GNN non temp.} & \textbf{One-Class SVM} & \textbf{Isol. Forest} \\
\hline
AUC aggregata                  & $\mathbf{0.965}\pm0.007$ & $0.555\pm0.037$ & $0.801\pm0.018$ & $0.763\pm0.014$ \\
AP aggregata                   & $\mathbf{0.922}\pm0.014$ & $0.560\pm0.021$ & $0.739\pm0.026$ & $0.575\pm0.016$ \\
\hline
Lateral --- AUC                & $\mathbf{0.913}\pm0.014$ & $0.451\pm0.021$ & $0.599\pm0.002$ & $0.645\pm0.025$ \\
Lateral --- Recall @1\%FPR     & $\mathbf{0.166}\pm0.009$ & $0.125\pm0.006$ & $0.062\pm0.016$ & $0.009\pm0.002$ \\
Recall aggregata @1\%FPR       & $\mathbf{0.665}\pm0.043$ & $0.366\pm0.004$ & $0.479\pm0.038$ & $0.116\pm0.020$ \\
\hline
\end{tabular}
\caption{Confronto sul test set sintetico (replay per-evento de-circolarizzato, \texttt{num\_events}=200\,000, \textbf{3 seed} \texttt{[42,7,123]}, media $\pm$ dev.std, soglia globale all'1\% FPR). Il vantaggio del TGN sul \textit{lateral movement} emerge isolando il segnale temporale: la GNN statica, priva di memoria ricorrente ma provvista degli stessi contatori tabellari, resta a livello di caso ($0.451\pm0.021$), contro l'AUC $0.913\pm0.014$ del TGN. Il recall laterale a soglia globale resta basso ($\sim$16.6\%) e l'apparente 12.5\% della GNN è \emph{spurio} (AUC$\approx$caso con soglia permissiva, non un segnale reale); la sua conversione in recall operativo è affidata alla calibrazione cost-sensitive instradata (Sez.~\ref{sec:threshold}). XGBoost \emph{supervisionato} (vede le etichette, fuori categoria) raggiunge AUC $0.974\pm0.001$ / lateral $0.929\pm0.001$: è un upper-bound di riferimento, non una baseline non supervisionata comparabile. \emph{Provenienza}: riga TGN da run \emph{v3} (keying per-cookie); le baseline sono rimisurate sotto la configurazione \emph{deployable} corrente (collasso \texttt{dev:guest}); \textit{Recall aggregata} alla soglia globale 1\% FPR su tutte le classi. Numeri rigenerati da \texttt{tests/regen\_report\_tables.py} (\texttt{tasks/runs/panelA.json}; cfr.\ \texttt{docs/latex/PROVENANCE.md}).}
\label{tab:baselines}
\end{table}

\subsubsection{Cost-sensitive routing}

Per convertire la forte capacità di ranking (lateral AUC $0.913\pm0.014$; Tab.~\ref{tab:baselines}) in recall operativo, è stato implementato un routing cost-sensitive. Il flusso di eventi viene separato a runtime in \textit{signal-dirty} e \textit{signal-clean}. Sul flusso pulito, la soglia viene rilassata per minimizzare i falsi negativi dei movimenti laterali (configurabile tramite un parametro di \texttt{cost\_ratio}). Sul test sintetico (3 seed, media), ciò porta il recall laterale dal $\sim$16.6\% (soglia globale all'1\% FPR, FPR benigno $\sim$1.6\%) al $\sim$35.0\% (decisione instradata, FPR benigno $\sim$5.0\%), con recall aggregato $\sim$75.8\%.

\section{L'evoluzione: il nodo \textit{Configuration} (schema v4)}
\label{sec:v4}

\subsection{Motivazione}

Il modello descritto nella Sezione~\ref{sec:v3} (schema v3) modella ogni richiesta come una catena causale a tre archi su quattro ruoli di nodo, \textit{source}$\rightarrow$\textit{device}$\rightarrow$\textit{user}$\rightarrow$\textit{resource}. In quello schema il \textit{fingerprint} TLS del client (JA3) è soltanto un \emph{bit di validità} (\texttt{features[0]}): segnala se l'handshake è regolare, ma non è un'\emph{identità}. Questo lascia scoperti due vettori d'attacco centrali nella ZTA:
\begin{itemize}
    \item \textit{Movimento laterale con tooling nuovo}: un dispositivo già noto, compromesso, inizia a parlare con un client/tool mai osservato su quella macchina. La validità TLS resta intatta (il tool usa una libreria regolare), ma il \emph{fingerprint} è nuovo.
    \item \textit{Furto di credenziali}: un attaccante che riusa le credenziali della vittima da un proprio client presenta un JA3 \emph{diverso} da quello abituale dell'utente, pur con ruolo/clearance leciti (l'attaccante eredita i permessi della vittima).
\end{itemize}
Entrambi sono \emph{signal-clean} e \emph{policy-clean}: indistinguibili sintatticamente dal traffico benigno se non per l'\emph{identità della configurazione}. Da qui la promozione della configurazione del client a \textbf{quinto ruolo di nodo}.

\subsection{Architettura del nodo \textit{Configuration}}

La configurazione è identificata da \texttt{conf:<ja3>} quando il \emph{fingerprint} è disponibile, altrimenti dal generico \texttt{conf:guest} (sempre presente lato serving, default server-side). Il JA3 mantiene quindi un \emph{doppio ruolo}, ortogonale: resta il bit di validità TLS in \texttt{features[0]} (catturato dalla \textit{Feature head}) \emph{e} diventa un'identità di nodo (catturata dalla memoria temporale e dalle \textit{history features}). Per costruzione la dimensione del messaggio (\texttt{msg\_dim}=7) e quella delle feature statiche di nodo (\texttt{node\_feat\_dim}=16) restano invariate: il nodo config, come i \textit{source}, non porta attributi statici dedicati ma solo memoria ricorrente, hashed-identity e slot di trust. Lo schema dell'evento passa a \texttt{schema\_version}=4 (i checkpoint v1/v2/v3 sono rifiutati al caricamento).

\subsubsection{Topologia: la catena causale a 5 archi}

Il nodo config è \emph{inserito tra source e device} nella catena causale, sostituendo l'arco diretto \textit{source}$\rightarrow$\textit{device} con la coppia \textit{source}$\rightarrow$\textit{config}$\rightarrow$\textit{device}, e aggiungendo un \emph{binding} discriminante \textit{config}$\rightarrow$\textit{user}. Il cammino più lungo è quindi:
\[
\textit{source} \;\rightarrow\; \textit{config} \;\rightarrow\; \textit{device} \;\rightarrow\; \textit{user} \;\rightarrow\; \textit{resource}.
\]

\begin{figure}[ht!]
\centering
\begin{tikzpicture}[
    font=\sffamily\small,
    nd/.style={draw, rounded corners, align=center, minimum height=0.85cm, minimum width=1.9cm, inner sep=3pt},
    src/.style={nd, fill=blue!8},
    cfg/.style={nd, fill=orange!18},
    dev/.style={nd, fill=blue!8},
    usr/.style={nd, fill=green!10},
    res/.style={nd, fill=red!8},
    ax/.style={->, >=Stealth, thick},
    bnd/.style={->, >=Stealth, thick, densely dashed},
    lbl/.style={font=\sffamily\scriptsize\itshape, midway}
]
\node[src] (S) {source\\\scriptsize IP};
\node[cfg, right=1.4cm of S] (C) {config\\\scriptsize JA3};
\node[dev, right=1.4cm of C] (D) {device\\\scriptsize TPM/cookie};
\node[usr, right=1.4cm of D] (U) {user};
\node[res, right=1.4cm of U] (R) {resource};

\draw[bnd] (S) -- (C) node[lbl, above]{bind};
\draw[bnd] (C) -- (D) node[lbl, above]{bind};
\draw[bnd] (D) -- (U) node[lbl, above]{bind};
\draw[ax]  (U) -- (R) node[lbl, above]{access};

% binding aggiuntivo config -> user
\draw[bnd] (C) to[out=-35, in=-145] node[lbl, below]{bind (furto cred.)} (U);
\end{tikzpicture}
\caption{Catena causale dello schema v4. I \textit{binding edge} (tratteggiati) portano messaggio-zero; il solo \textit{access edge} \textit{user}$\rightarrow$\textit{resource} porta il messaggio a 7 dimensioni. Lo score dell'evento è il \emph{massimo} sui suoi archi. L'arco diretto \textit{source}$\rightarrow$\textit{device} di v3 è sostituito da \textit{source}$\rightarrow$\textit{config}$\rightarrow$\textit{device}; il \textit{binding} \textit{config}$\rightarrow$\textit{user} discrimina il furto di credenziali (un client diverso da quello abituale della vittima).}
\label{fig:v4chain}
\end{figure}

\begin{table}[ht!]
\centering
\renewcommand{\arraystretch}{1.3}
\begin{tabular}{l l l}
\hline
\textbf{Arco} & \textbf{Tipo} & \textbf{Stato} \\
\hline
\textit{user} $\rightarrow$ \textit{resource}  & access (msg 7-dim) & esistente \\
\textit{device} $\rightarrow$ \textit{user}    & binding (msg-zero) & esistente \\
\textit{source} $\rightarrow$ \textit{config}  & binding (msg-zero) & nuovo (v4) \\
\textit{config} $\rightarrow$ \textit{device}  & binding (msg-zero) & nuovo (v4) \\
\textit{config} $\rightarrow$ \textit{user}    & binding (msg-zero) & nuovo (v4) \\
\hline
\end{tabular}
\caption{Edge set per richiesta nello schema v4 (config sempre presente lato serving). Ordine di commit in memoria, identico in training/replay/serving: \textit{source}$\rightarrow$\textit{config}, \textit{config}$\rightarrow$\textit{device}, \textit{config}$\rightarrow$\textit{user}, \textit{device}$\rightarrow$\textit{user}, \textit{user}$\rightarrow$\textit{resource}. Fallback coerenti con la gestione opzionale già esistente: device assente $\Rightarrow$ \textit{config}$\rightarrow$\textit{user} fa da ponte; source assente $\Rightarrow$ si salta solo \textit{source}$\rightarrow$\textit{config}; per dataset senza config (es. LANL) il \textit{gating} \texttt{has\_config} ripristina il percorso legacy \textit{source}$\rightarrow$\textit{device}.}
\label{tab:v4edges}
\end{table}

L'architettura core (memoria ricorrente, \textit{GraphAttentionEmbedding} multi-hop, \textit{Dual Head Scorer}) è \emph{invariata}: il modello è agnostico al numero di archi per evento, e i nuovi \textit{binding} riusano la stessa primitiva di scoring degli archi esistenti (con \textit{history features} senza ausiliario, come per \textit{device}$\rightarrow$\textit{user}). Il generatore sintetico è esteso coerentemente: ogni macchina riceve 1--2 configurazioni abituali da un pool condiviso; il movimento laterale presenta con probabilità $0.5$ un \emph{tool} (config) mai usato da quel device; il furto di credenziali alloca al client attaccante un JA3 mai visto (\texttt{conf:atk-NNN}). \emph{Inoltre, per supportare rotte pubbliche (es. \texttt{/login}) che non richiedono un JWT e prevenire falsi positivi da cold-start, lo stream inietta regolarmente accessi a tali risorse usando lo slot utente \texttt{anonymous} combinato al fallback \texttt{dev:guest}, insegnando al modello la normalità del traffico non autenticato.}

% ============================ RISULTATI v4 ============================
\subsection{Risultati: il contributo del nodo config}
\label{sec:v4results}

La valutazione segue lo stesso protocollo della Sezione~\ref{sec:v3} (replay per-evento attraverso l'esatto codice di serving, split cronologico 70/10/20, soglia cost-sensitive instradata). Si riportano due confronti complementari: (i) il modello v3 vs v4 (\textit{keying} per-cookie) all'\textit{operating point} instradato (\textbf{3 seed} \texttt{[42,7,123]}, media $\pm$ dev.std, 200\,000 eventi / 15 epoche), e (ii) la \emph{validazione mirata} del nodo config — un'ablazione controllata \emph{a parità di dati} su seed multipli.

\subsubsection{Confronto v3 $\rightarrow$ v4 (\textit{keying} per-cookie)}

La Tabella~\ref{tab:v3v4} confronta l'artefatto (\textit{keying} per-cookie, coerente con la riga TGN della Tab.~\ref{tab:baselines}) prima e dopo l'introduzione del nodo config, al medesimo \textit{operating point} (decisione cost-sensitive instradata sul segnale), su 3 seed. Il guadagno \emph{robusto} è sul \textbf{recall operativo del movimento laterale}: passa da $\sim$35.0\% a $\sim$57.7\% ($+0.227$, oltre la banda di rumore multi-seed), coerente con la validazione multi-seed \textit{theft-rich} (Tab.~\ref{tab:theft}, $+0.135$). A differenza della precedente valutazione single-run, sui 3 seed anche l'\textbf{AUC laterale} migliora in modo significativo ($0.913\pm0.014\rightarrow0.930\pm0.013$, $+0.017$), mentre le metriche aggregate restano \emph{piatte entro il rumore} (AUC aggregata $\approx$invariata $0.965\rightarrow0.964$; AP aggregata $0.922\rightarrow0.897$, $-0.025$ non significativo) e l'FPR benigno instradato \emph{sale} ($\sim$5.0\% $\rightarrow$ $\sim$6.6\%, $+0.017$): il nodo config non domina uniformemente, ma \emph{sposta} il modello verso un recall operativo del laterale più alto al punto di lavoro instradato --- il \textit{trade-off} che si paga è un FPR benigno maggiore e un'AP aggregata leggermente più bassa. È il comportamento atteso per costruzione: \textit{config}$\rightarrow$\textit{device} cattura il ``tool nuovo su device noto'' e \textit{config}$\rightarrow$\textit{user} il ``client che quell'utente non usa abitualmente'', cioè i \textit{tell} del laterale e del furto di credenziali.

\begin{table}[ht!]
\centering
\footnotesize
\renewcommand{\arraystretch}{1.4}
\setlength{\tabcolsep}{5pt}
\begin{tabular}{l c c c}
\hline
\textbf{Metrica (test set)} & \textbf{v3 (4 nodi)} & \textbf{v4 (nodo config)} & $\Delta$ \\
\hline
AUC aggregata                    & $\mathbf{0.965}\pm0.007$ & $0.964\pm0.010$ & $\approx\!0^{\dagger}$ \\
AP aggregata                     & $\mathbf{0.922}\pm0.014$ & $0.897\pm0.038$ & $-0.025^{\dagger}$ \\
Recall aggregata (instradata)    & $0.758\pm0.030$ & $\mathbf{0.829}\pm0.040$ & $+0.071$ \\
\hline
Lateral --- AUC                  & $0.913\pm0.014$ & $\mathbf{0.930}\pm0.013$ & $+0.017$ \\
Lateral --- AP                   & $0.522\pm0.020$ & $\mathbf{0.575}\pm0.098$ & $+0.052^{\dagger}$ \\
Lateral --- Recall (instradata)  & $0.350\pm0.034$ & $\mathbf{0.577}\pm0.057$ & $+0.227$ \\
Lateral --- Recall (@1\%FPR)     & $0.166\pm0.009$ & $\mathbf{0.363}\pm0.128$ & $+0.197$ \\
\hline
Policy --- AUC                   & $\mathbf{0.987}\pm0.007$ & $0.972\pm0.015$ & $-0.015$ \\
Contextual --- AUC               & $\mathbf{0.994}\pm0.001$ & $0.991\pm0.002$ & $-0.003$ \\
\hline
FPR benigno (instradato)         & $\mathbf{0.050}\pm0.002$ & $0.066\pm0.013$ & $+0.017$ \\
\hline
\end{tabular}
\caption{Confronto v3 vs v4 (\textit{keying} per-cookie) sul medesimo test set sintetico (replay per-evento, \texttt{num\_events}=200\,000, \textbf{3 seed} \texttt{[42,7,123]}, media $\pm$ dev.std, 15 epoche, decisione cost-sensitive instradata, \texttt{save=False}). Il marcatore $^{\dagger}$ indica un $\Delta$ \emph{entro} la banda di rumore multi-seed ($|\Delta|\le$ max delle due dev.std), quindi non significativo: AUC aggregata ($\approx\!0$) e AP aggregata ($-0.025$). I guadagni \emph{robusti} del nodo config sono il recall operativo del laterale instradato ($+0.227$), il recall aggregato instradato ($+0.071$) e --- a differenza della precedente valutazione single-run --- l'AUC laterale ($+0.017$, ora oltre il rumore), al costo di un FPR benigno più alto ($+0.017$). Il recall laterale @1\%FPR ha varianza inter-seed elevata per v4 ($\pm0.128$) e va letto con cautela. Le classi \textit{policy} (OPA-owned) e \textit{contextual} (rule-trivial) restano sostanzialmente invariate (recall@thr $\sim$0.93--0.98 per entrambe). Il confronto con le baseline non-relazionali (Tab.~\ref{tab:baselines}) è immutato: nessuna baseline modella il nodo config. Numeri rigenerati da \texttt{tests/regen\_report\_tables.py} (\texttt{tasks/runs/panelB.json}).}
\label{tab:v3v4}
\end{table}

\subsubsection{Validazione mirata: il furto di credenziali}

L'\textit{operating point} pubblicato non permetteva di misurare il furto di credenziali: gli slot attaccante del generatore si esauriscono presto (configurazione di default \texttt{num\_theft\_slots}=64, \texttt{p\_cred\_theft}=0.0012 $\Rightarrow$ tutti gli incidenti partono entro $\sim$l'evento 53k, dentro il 70\% di \emph{train}), lasciando \emph{zero} incidenti di furto e \emph{zero} cookie-wipe nello split di test. Per quantificare il contributo specifico del binding \textit{config}$\rightarrow$\textit{user} si è quindi eseguita una \textbf{validazione mirata} su uno stream \emph{theft-rich} (slot e tasso d'incidente aumentati \emph{solo per questa valutazione}, così che gli incidenti si distribuiscano sull'intero stream incluso il test), confrontando — \emph{a parità di dati}, su 3 \texttt{seed} (media$\pm$dev.std) e con \texttt{save=False} (l'artefatto in \texttt{public/} non viene toccato) — il modello v4 completo contro la sua \emph{ablazione del nodo config} (il modello ricade sul percorso legacy \textit{source}$\rightarrow$\textit{device}, cioè $\approx$v3 sugli stessi dati). I risultati sono in Tabella~\ref{tab:theft}.

\begin{table}[ht!]
\centering
\renewcommand{\arraystretch}{1.4}
\begin{tabular}{l c c c}
\hline
\textbf{Metrica (test, theft-rich)} & \textbf{v4 (nodo config)} & \textbf{$\approx$v3 (ablazione)} & $\Delta$ \\
\hline
Furto credenziali --- AUC        & $0.729\pm0.088$ & $0.676\pm0.036$ & $+0.053$ \\
Furto credenziali --- Recall@thr & $\mathbf{0.314}\pm0.107$ & $0.206\pm0.039$ & $+0.108$ \\
\hline
Lateral --- AUC                  & $0.919\pm0.018$ & $0.913\pm0.007$ & $+0.006$ \\
Lateral --- Recall@thr           & $\mathbf{0.469}\pm0.088$ & $0.334\pm0.014$ & $+0.135$ \\
\hline
AUC aggregata                    & $0.958\pm0.011$ & $0.956\pm0.005$ & $+0.002$ \\
Recall aggregata (instradata)    & $0.777\pm0.047$ & $0.744\pm0.005$ & $+0.033$ \\
FPR benigno (cookie-wipe)        & $0.060\pm0.021$ & $0.038\pm0.033$ & $+0.022$ \\
\hline
\end{tabular}
\caption{Validazione mirata del nodo config (3 \texttt{seed} [42,7,123], 80\,000 eventi / 12 epoche, stream \emph{theft-rich}: slot 400/120, \texttt{p\_theft}=0.002, \texttt{p\_wipe}=0.0008; \texttt{save=False}). Lo split mirato porta finalmente $n=452$ eventi di furto nel test (erano $0$). L'ablazione disattiva il nodo config \emph{a parità di stream} (percorso legacy \textit{source}$\rightarrow$\textit{device}). Il nodo config aumenta il \textit{recall} sul furto credenziali ($+0.108$) e sul laterale ($+0.135$); l'AUC sul furto migliora ($+0.053$) ma con dev.std ampia (conteggi per-seed modesti), quindi il segnale robusto è il \textit{recall} e il fatto che la detection diventi \emph{possibile}. L'aggregata è invariata. Si osserva un lieve aumento del FPR benigno sui cookie-wipe ($+0.022$, entro le bande di dev.std): il binding \textit{config}$\rightarrow$\textit{user} è leggermente più sensibile sui dispositivi re-keyed, trade-off accettabile dato il guadagno sul recall.}
\label{tab:theft}
\end{table}

\subsubsection{Sweep architetturale}
\label{sec:archsweep}

L'introduzione di un quinto nodo (e di due archi aggiuntivi per richiesta) suggerisce di verificare se al modello serva \emph{più capacità}. Si sono valutate tre direzioni — un layer nascosto aggiuntivo nell'MLP della \textit{Feature head} (\texttt{link\_pred\_hidden\_layers}=3), una memoria ricorrente più ampia (\texttt{memory\_dim}=384) e più teste di attenzione nei \textit{TransformerConv} (\texttt{gnn\_heads}=8) — più la loro combinazione, sempre su 3 \texttt{seed} con \texttt{save=False}. La Tabella~\ref{tab:archsweep} riporta l'esito.

\begin{table}[ht!]
\centering
\footnotesize
\renewcommand{\arraystretch}{1.4}
\setlength{\tabcolsep}{5pt}
\begin{tabular}{l c c c c c}
\hline
\textbf{Variante} & \textbf{Lateral AUC} & \textbf{Lateral AP} & \textbf{Lateral Rec@thr} & \textbf{agg AUC} & \textbf{agg Recall} \\
\hline
\textbf{baseline v4}     & $0.929\pm0.009$ & $\mathbf{0.645}\pm0.023$ & $\mathbf{0.607}\pm0.060$ & $\mathbf{0.966}\pm0.005$ & $\mathbf{0.832}\pm0.019$ \\
$+1$ layer MLP           & $0.931\pm0.004$ & $0.641\pm0.035$ & $0.560\pm0.037$ & $0.965\pm0.005$ & $0.803\pm0.029$ \\
$+$memory (384)          & $0.905\pm0.015$ & $0.543\pm0.051$ & $0.441\pm0.050$ & $0.957\pm0.006$ & $0.766\pm0.015$ \\
$+$heads (8)             & $0.896\pm0.037$ & $0.558\pm0.130$ & $0.459\pm0.161$ & $0.956\pm0.015$ & $0.777\pm0.061$ \\
$+$mem$+$heads$+$layer    & $0.913\pm0.025$ & $0.526\pm0.105$ & $0.443\pm0.111$ & $0.958\pm0.011$ & $0.763\pm0.039$ \\
\hline
\end{tabular}
\caption{Sweep architetturale (3 \texttt{seed} [42,7,123], 40\,000 eventi / 12 epoche, \texttt{save=False}). Nessun aumento di capacità migliora la classe \textit{model-owned} (movimento laterale): il layer MLP extra è \emph{neutro} sull'AUC ($+0.002$, entro la dev.std) ma riduce il \textit{recall} operativo; memoria più ampia e più teste \emph{peggiorano} AUC/AP/recall e \emph{aumentano} la varianza tra seed (es. \texttt{gnn\_heads}=8: recall $0.459\pm0.161$). Il collo di bottiglia del laterale non è la capacità parametrica ma il \emph{segnale} (history features + memoria temporale), coerentemente con l'ablation a 3 seed della Sez.~\ref{sec:v3} (history features $+0.163$ AUC, dominante). \textbf{Decisione}: si mantiene l'architettura corrente; il \textit{deployable} non viene modificato.}
\label{tab:archsweep}
\end{table}

\subsubsection{Sintesi}

Il nodo \textit{configuration} (v4) alza sostanzialmente il \textbf{recall operativo} del movimento laterale al punto di lavoro instradato ($+0.227$ a 3 seed) e abilita la rilevazione del furto di credenziali (recall $+0.108$ vs l'ablazione a parità di dati), senza costi architetturali: \texttt{msg\_dim} e \texttt{node\_feat\_dim} restano invariati, e lo sweep conferma che la capacità attuale è già adeguata. Non è però un dominio uniforme: il guadagno è \emph{sul recall instradato} (e sull'AUC laterale, $+0.017$, ora oltre il rumore), e si paga con un FPR benigno più alto ($+0.017$) e un'AP aggregata leggermente più bassa ($-0.025$, entro il rumore), mentre l'AUC aggregata resta piatta entro il rumore multi-seed; è quindi uno \textit{spostamento dell'operating point} verso il laterale, non un miglioramento su tutte le metriche. Questo guadagno operativo si \emph{somma} a quello già documentato del TGN sulle baseline non-relazionali (Tab.~\ref{tab:baselines}), che non modellano la configurazione del client.

% ==================== ESPERIMENTO: identità del device ====================
\subsection{Esperimento: granularità dell'identità del dispositivo (nodo \texttt{dev:guest})}
\label{sec:guestdev}

\subsubsection{Ipotesi e motivazione}

Nello schema v4 il nodo \textit{device} è identificato da \texttt{tpm:<id>} quando il TPM è attestato e, in assenza di TPM, da un \emph{cookie persistente per-macchina} (\texttt{ck:<id>}), con re-keying su un nodo ``freddo'' a ogni \textit{cookie-wipe}. Ogni macchina senza TPM è quindi un \emph{nodo distinto} del grafo. In analogia con il nodo \textit{configuration} --- dove un \textit{fingerprint} non risolvibile ricade sempre sul generico \texttt{conf:guest} (Sez.~\ref{sec:v4}) --- ci si è chiesti se convenga \emph{collassare tutti i dispositivi senza TPM su un unico nodo condiviso} \texttt{dev:guest}, anonimo e a bassa fiducia, invece di mantenere un'identità-cookie per macchina. L'ipotesi indaga un \emph{trade-off di modellazione}: l'identità per-cookie aggiunge potere discriminante (un device mai visto è una \textit{novelty cue}) oppure è soltanto \emph{rumore}, dato che gli attacchi \textit{rule-blind} (laterale, furto credenziali) sono per costruzione \textit{signal-clean} e il loro segnale è già portato dalle edge \textit{config}/\textit{source}/\textit{user}?

Il comportamento è esposto da un flag di configurazione (\texttt{guest\_device\_fallback}, default \emph{on} = collasso su \texttt{dev:guest}, adottato come configurazione deployabile; disattivando il flag si ripristina il keying per-cookie per-macchina): quando attivo, ogni device non-TPM (cookie, IP-debole o assente) collassa su \texttt{dev:guest} sia nel generatore (instradamento sul medesimo nodo, poiché in addestramento i nodi sono indici di \textit{slot}, non chiavi) sia lato serving (collasso della chiave esterna), così che addestramento e \textit{serving} restino coerenti. Coerentemente, il \textit{cookie-wipe} viene \emph{neutralizzato}: un nodo guest condiviso non possiede un cookie per-macchina da resettare. È importante notare che, sotto questa politica, \emph{anche il dispositivo dell'attaccante} nel furto di credenziali (privo di TPM) collassa su \texttt{dev:guest}: il \textit{tell} di ``device mai visto'' sparisce e la rilevazione deve appoggiarsi unicamente all'IP e al JA3 mai visti e al \textit{binding} \textit{config}$\rightarrow$\textit{user}.

\subsubsection{Protocollo e risultati}

Si confronta --- a parità di impostazioni del generatore, su 3 \texttt{seed} [42,7,123] (media$\pm$dev.std), con \texttt{save=False} e sullo stesso stream \emph{theft-rich} della validazione del nodo config (80\,000 eventi / 12 epoche, slot 400/120, \texttt{p\_theft}=0.002, \texttt{p\_wipe}=0.0008) --- il \textit{keying} per-cookie (\emph{baseline}) contro il collasso su \texttt{dev:guest}. La Tabella~\ref{tab:guestdev} riporta l'esito.

\begin{table}[ht!]
\centering
\renewcommand{\arraystretch}{1.4}
\begin{tabular}{l c c c}
\hline
\textbf{Metrica (test, theft-rich)} & \textbf{cookie (baseline)} & \textbf{\texttt{dev:guest}} & $\Delta$ \\
\hline
Lateral --- AUC                  & $0.915\pm0.019$ & $0.911\pm0.007$ & $-0.003$ \\
Lateral --- AP                   & $0.583\pm0.068$ & $\mathbf{0.632}\pm0.017$ & $+0.049$ \\
Lateral --- Recall@thr           & $0.453\pm0.060$ & $\mathbf{0.488}\pm0.019$ & $+0.035$ \\
\hline
Furto credenziali --- AUC        & $0.701\pm0.078$ & $\mathbf{0.804}\pm0.024$ & $+0.103$ \\
Furto credenziali --- Recall@thr & $0.244\pm0.082$ & $\mathbf{0.287}\pm0.062$ & $+0.043$ \\
\hline
AUC aggregata                    & $0.954\pm0.013$ & $\mathbf{0.961}\pm0.001$ & $+0.007$ \\
Recall aggregata (instradata)    & $0.764\pm0.036$ & $\mathbf{0.787}\pm0.006$ & $+0.023$ \\
FPR benigno (instradato)         & $0.045\pm0.002$ & $\mathbf{0.037}\pm0.004$ & $-0.007$ \\
FPR cookie-wipe                  & $0.041\pm0.013$ ($n{=}651$) & --- ($n{=}0$) & neutralizzato \\
\hline
\end{tabular}
\caption{Esperimento sull'identità del device: \textit{keying} per-cookie vs collasso su \texttt{dev:guest} (3 \texttt{seed} [42,7,123], 80\,000 eventi / 12 epoche, stream \emph{theft-rich}, \texttt{save=False}). Le due varianti condividono \texttt{seed} e impostazioni ma \emph{non} sono appaiate evento-per-evento: la neutralizzazione del \textit{cookie-wipe} altera il percorso del generatore di numeri casuali (da cui il diverso conteggio di positivi: \textit{lateral} $n{=}3959$ vs $4032$, furto $n{=}452$ vs $386$ sui 3 \texttt{seed}); il confronto è quindi \emph{distribuzionale}. I valori $\Delta$ sono calcolati sulle medie non arrotondate (donde l'apparente $-0.003$ a fronte di $0.915$ vs $0.911$ in tabella). Il collasso su \texttt{dev:guest} non degrada la detection (lateral AUC piatto, $-0.003$ entro la dev.std) e anzi migliora il furto credenziali, abbassa l'FPR benigno, \emph{riduce marcatamente la varianza tra seed} ed elimina l'intera classe di falsi positivi da cookie-wipe ($n_{\text{wiped}}{=}0$, conferma della neutralizzazione implementata).}
\label{tab:guestdev}
\end{table}

\subsubsection{Interpretazione}

Il risultato è \emph{controintuitivo}: rimuovere l'identità per-macchina dei dispositivi senza TPM non solo non danneggia la rilevazione, ma costituisce un \emph{miglioramento di Pareto} sulle metriche misurate. La lettura è coerente con l'architettura del modello:
\begin{itemize}
    \item \textbf{L'identità per-cookie era rumore netto.} I nodi-cookie per-macchina sono \emph{sparsi}: ciascuno vede pochi eventi, quindi accumula una memoria ricorrente e dei contatori di storia deboli e freddi. Concentrare tutto il traffico non-TPM su un solo nodo \texttt{dev:guest} produce \textit{binding} \textit{config}$\rightarrow$\textit{device} e \textit{device}$\rightarrow$\textit{user} statisticamente più stabili, riducendo i falsi positivi benigni ($-0.007$ FPR) e --- in modo netto --- la varianza tra \texttt{seed} (es. lateral AUC $\pm0.007$ vs $\pm0.019$; furto AUC $\pm0.024$ vs $\pm0.078$).
    \item \textbf{Il segnale degli attacchi non risiedeva nel device.} Laterale e furto credenziali sono \textit{signal-clean} per costruzione: la \textit{novelty cue} vive sulle \textit{history features} delle edge \textit{user}/\textit{config}/\textit{source}, non sulla novità del nodo-device. Per il furto, pur perdendo il \textit{tell} di ``device mai visto'' (l'attaccante collassa anch'esso su \texttt{dev:guest}), restano l'IP e il JA3 mai visti e il \textit{binding} \textit{config}$\rightarrow$\textit{user} (già documentato come discriminante principale del furto, Sez.~\ref{sec:v4results}); il netto guadagno di AUC del furto ($+0.103$) riflette soprattutto la \emph{riduzione di varianza} rispetto a una \textit{baseline} ad alta dev.std ($\pm0.078$).
    \item \textbf{Eliminazione di una classe di falsi positivi.} Senza identità-cookie non esiste il \textit{re-keying} su nodo freddo: l'intera classe di falsi positivi da \textit{cookie-wipe} (FPR $\sim$0.041 su $n{=}651$ nella \textit{baseline}) scompare ($n_{\text{wiped}}{=}0$).
\end{itemize}

\paragraph{Limiti e costo non misurato.} Il guadagno sulle metriche di detection va soppesato con un costo che esse \emph{non} catturano: collassando i dispositivi senza TPM si perde l'\emph{attribuzione per-macchina} (tracciabilità forense, capacità di isolare un singolo endpoint compromesso, audit). Inoltre il confronto è distribuzionale e su stream sintetico: il $\Delta$ sull'AUC laterale è entro il rumore e la quota di guadagno dovuta alla riduzione di varianza non è separabile da un eventuale guadagno ``vero''. Nondimeno, dato il miglioramento di Pareto misurato (FPR benigno e varianza tra seed più bassi, eliminazione dei falsi positivi da cookie-wipe), il \textit{deployable} ha \emph{adottato} il collasso su \texttt{dev:guest} come \textit{default} (\texttt{guest\_device\_fallback=True}); il keying per-cookie resta una politica attivabile (disabilitando il flag) per i contesti che richiedono l'\emph{attribuzione forense per-macchina} --- il costo consapevolmente accettato di questa scelta. Si noti che le Tab.~\ref{tab:v3v4} e~\ref{tab:baselines} descrivono l'artefatto deployabile \emph{multi-seed} (3 seed \texttt{[42,7,123]}) col keying per-cookie, prodotto \emph{prima} di questa adozione; i loro valori restano validi come tali. L'evidenza qui resta inoltre utile a mostrare che, su questo stream, l'identità fine del dispositivo \emph{non} è il portatore del segnale.

% ==================== DEPLOYMENT ====================
\section{Deployment e specifiche hardware/software}
\label{sec:deploy}

Le primitive di model-serving (estrazione memoria ed espansione del grafo storico) operano asintoticamente a tempo costante $O(1)$, e il TGN mantiene tutto il suo stato dinamico e ricorrente nativamente in RAM. Da questa scelta discendono due conseguenze di deployment: il servizio è \emph{single-worker} --- lo stato comportamentale per-entità non è replicabile senza frammentarlo (limite discusso in Sez.~\ref{sec:limiti}) --- ed è fondamentale il dimensionamento asincrono.

\subsection{Piattaforma hardware/software (addestramento e inferenza)}

Lo stesso container isolato ospita sia l'addestramento sia l'inferenza. Al fine di calcolare la backward-propagation e smaltire la mole della \textit{Graph Attention}, esso si appoggia alle seguenti specifiche hardware/software:

\begin{itemize}
    \item container engine nativo agganciato a librerie NVIDIA CUDA 13;
    \item hardware accelerato su GPU con architettura NVIDIA Blackwell (testato su RTX 5070 Ti). La GPU accelera la \textit{Graph Attention} multi-hop in inferenza e la backward-propagation in addestramento; lo stato ricorrente della \texttt{TGNMemory} e i ring-buffer del \texttt{MessageNeighborLoader} sono mantenuti con accesso $O(1)$ su tensori preallocati.
\end{itemize}

Nel momento di chiusura del framework (o tramite specifici endpoint API persistenti), la memoria serializzata di PyTorch (\texttt{tgn\_checkpoint.pt} e \texttt{tgn\_stats.json}) va copiata su sistema persistente per evitare cold-start traumatici.

\subsection{Serving e latenza di inferenza}

Il sistema finale eroga un'infrastruttura single-worker fondata su \texttt{FastAPI} (gestione I/O asincrono per polling concorrenziale) che elabora gli eventi con una latenza end-to-end (round-trip HTTP client$\to$servizio sulla rete container, catena v4 completa a 5 archi \textit{source}$\to$\textit{config}$\to$\textit{device}$\to$\textit{user}$\to$\textit{resource}) di $P50 \approx 9.6$ ms ($P99 \approx 10.8$ ms); l'omissione del binding di dispositivo non altera apprezzabilmente né la latenza né la detection in cold-start. Misure su GPU RTX 5070 Ti (VRAM occupata $\sim$2.5\,GB), tramite test client di streaming live in regime \textit{cold-start} (entità mai viste prima, prefisso \texttt{prod\_}): un sanity-check d'induttività più che un punto operativo, poiché in cold-start il modello opera ad alto recall ma bassa specificità (lateral recall $\sim$82\%, specificità benigna $\sim$43\% --- FPR gonfiato dall'assenza di storico). \emph{Le feature di nodo sono ora applicate per-tipo (tier del dispositivo solo sul nodo \textit{device}, attributi nulli sul nodo \textit{utente}, coerentemente col contratto di training): correggere il precedente \textit{train/serve skew} --- in cui il tier veniva scritto anche sul nodo utente, sempre nullo in training --- riduce il falso positivo da cold-start (specificità benigna da $\sim$33\% a $\sim$43\% sul medesimo checkpoint, a fronte di un recall laterale da $\sim$87\% a $\sim$82\%).} In puro cold-start il binding di dispositivo è trascurabile (il device, anch'esso privo di storia, non porta segnale aggiuntivo); il punto operativo affidabile resta quello calibrato offline (Sez.~\ref{sec:eval}), non influenzato da questa correzione (il replay non inietta le feature statiche per-evento).

% ==================== LIMITI ====================
\section{Limiti del modello}
\label{sec:limiti}

% \subsection{Il paradosso dell'induttività (cold start)}

% Affinché il TGN riconosca deviazioni strutturali di alto livello, la rete neurale deve esplorare la baseline (\textit{Temporal Neighborhood}) su un orizzonte storico. Al primo accesso, un utente appena iscritto risulterebbe sempre in uno stato di fallimento topologico. Durante questa prima fase (cold start), l'Orchestrator OPA è sprovvisto dell'intelligenza del sistema di Intelligenza Artificiale. Per abbassare tale superficie d'attacco, si limita il buio di validazione con un \textit{Grace Period} ristretto a soli 5 eventi.

\subsection{Assenza di scalabilità orizzontale}

Il limite di deployment più strutturale discende direttamente dal design illustrato in Sez.~\ref{sec:deploy}: lo stato comportamentale per-entità --- la memoria ricorrente della \texttt{TGNMemory}, i ring-buffer dei vicini temporali del \texttt{MessageNeighborLoader} e i contatori espliciti di \textit{history} --- risiede interamente in RAM all'interno di un \emph{singolo worker}. Replicare il servizio orizzontalmente (più repliche dietro un \textit{load balancer}) frammenterebbe questo database comportamentale: ogni replica osserverebbe soltanto una frazione dello stream e accumulerebbe una memoria e dei contatori \emph{parziali e incoerenti}, degradando proprio il segnale temporale su cui poggia la rilevazione del movimento laterale. Non esiste inoltre uno stato condiviso a bassa latenza consultabile in tempo $O(1)$ a ogni evento senza reintrodurre quel database a grafo che l'architettura ha esplicitamente eliminato. La scalabilità del sistema è quindi \emph{verticale} (GPU e RAM più capienti, per sostenere grafi e memorie più ampi) e \emph{asincrona} (l'I/O non bloccante di \texttt{FastAPI} assorbe il carico concorrenziale su un singolo processo), non orizzontale. Una distribuzione del traffico tra repliche richiederebbe uno \textit{sharding} coerente per-entità (ogni entità sempre instradata alla medesima replica), con la complessità e i punti di fallimento che ciò comporta --- una direzione di lavoro futura, non supportata dall'artefatto corrente.

\subsection{La sfida del lateral movement}

Senza l'ausilio delle etichette (unsupervised), il lateral movement non mostra tracce o indicatori palesi: è feature-identico a un accesso benigno non abituale. Le leve che lo rendono \emph{ordinabile} (lateral AUC $0.913\pm0.014$ a 3 seed a piena capacità, keying per-cookie; da $\sim$0.45 della GNN statica) sono il vicinato temporale ampio ($\text{num\_hops}{=}3$, $\text{neighbor\_size}{=}30$, $\text{memory\_dim}{=}256$) e soprattutto le \textit{history features} esplicite ($+0.163$ AUC, il contributo dominante), affiancate dall'hashed-identity ($+0.046$ AUC); il precursore di kill-chain ($+0.013$ AUC) e la testa strutturale ($+0.007$ AUC) risultano invece \emph{marginali} (ablation multi-seed). La conversione del ranking in recall operativo avviene poi con la soglia cost-sensitive instradata (Sez.~\ref{sec:threshold}): sul test sintetico il recall laterale passa da $\sim$\textit{16.6\%} (soglia globale all'1\% FPR, FPR benigno $\sim$1.6\%) a $\sim$\textit{35.0\%} (decisione instradata, FPR benigno $\sim$5.0\%) --- operating point regolabile via \texttt{cost\_ratio}/\texttt{clean\_fpr\_cap}. Resta la classe più ostica (AP $\sim$0.52), ma nettamente sopra le baseline (cfr. Sez.~\ref{sec:eval}).

\subsection{Anti-poisoning gate nel live stream}

Il design offline del testing replica 1:1 il routing live, prevenendo fenomeni di \textit{data leakage}. Tuttavia, l'ambiente di deployment espone alla minaccia di \textit{data poisoning}. Un aggressore persistente potrebbe infiltrarsi immettendo attacchi lenti per abituare il grafo ai propri comportamenti nocivi, sfalsando lentamente lo scoring e la similarità del coseno. Il problema è mitigato dal gate anti-avvelenamento: sia l'internal state della \texttt{TGNMemory} che le code a scorrimento dei vicini temporali vengono aggiornate e \textit{committate} in RAM solo nel caso in cui lo score della rete si posizioni sotto il limite calcolato per le minacce ($< threshold$). I pacchetti classificati come anomalia, inviati a OPA, vengono ignorati dal sistema di Intelligenza Artificiale, preservando incontaminato lo stato vettoriale del modello.

% ==================== QUADRO RIEPILOGATIVO ====================
\clearpage
\section{Quadro riepilogativo: confronto di tutti i modelli}
\label{sec:summary}

Le valutazioni discusse nelle due sezioni-modello sono state condotte sotto protocolli diversi (stream sintetico standard vs \emph{theft-rich}, soglia globale all'1\% FPR vs decisione cost-sensitive instradata; tutte multi-seed su 3 seed) e non sono quindi confrontabili \emph{tra} di loro in valore assoluto. La Tabella~\ref{tab:summary} le consolida in un unico quadro, organizzato in tre pannelli \emph{per protocollo}: \textbf{sono validi solo i confronti all'interno dello stesso pannello}. Lo scopo è dare una lettura d'insieme della progressione --- dalle baseline non relazionali al TGN, e dal modello a 4 nodi (v3) a quello con nodo \textit{configuration} (v4) --- senza dover ricomporre i numeri sparsi nelle Tabelle~\ref{tab:baselines},~\ref{tab:v3v4},~\ref{tab:theft} e~\ref{tab:guestdev}.

\begin{table}[ht!]
\centering
\footnotesize
\renewcommand{\arraystretch}{1.3}
\setlength{\tabcolsep}{5pt}
\begin{tabular}{l c c c c c c}
\toprule
\textbf{Modello / variante} & \textbf{AUC agg} & \textbf{AP agg} & \textbf{Lat.\ AUC} & \textbf{Lat.\ AP} & \textbf{Lat.\ Recall} & \textbf{agg Recall} \\
\midrule
\multicolumn{7}{l}{\textit{Pannello A — stream standard, 3 seed [42,7,123] (media, 200k ev.), soglia globale 1\% FPR}} \\
\addlinespace[2pt]
\textbf{TGN} (v3, per-cookie)            & \textbf{0.965} & \textbf{0.922} & \textbf{0.913} & ---   & 0.166 & \textbf{0.665} \\
GNN non temporale (ablation)             & 0.555 & 0.560 & 0.451 & ---   & 0.125$^{\dagger}$ & 0.366 \\
One-Class SVM                            & 0.801 & 0.739 & 0.599 & ---   & 0.062 & 0.479 \\
Isolation Forest                         & 0.763 & 0.575 & 0.645 & ---   & 0.009 & 0.116 \\
XGBoost (supervisionato)$^{\ddagger}$    & 0.974 & 0.948 & 0.929 & ---   & 0.339 & 0.772 \\
\midrule
\multicolumn{7}{l}{\textit{Pannello B — stream standard, 3 seed [42,7,123] (media), decisione cost-sensitive instradata}} \\
\addlinespace[2pt]
v3 (4 nodi)                              & \textbf{0.965} & \textbf{0.922} & 0.913 & 0.522 & 0.350 & 0.758 \\
\textbf{v4} (nodo config)                & 0.964 & 0.897 & \textbf{0.930} & \textbf{0.575} & \textbf{0.577} & \textbf{0.829} \\
\midrule
\multicolumn{7}{l}{\textit{Pannello C — stream theft-rich, 3 seed (media), ablazioni controllate (\texttt{save=False})}} \\
\addlinespace[2pt]
v4 (nodo config)                         & 0.958 & ---   & 0.919 & ---   & 0.469 & 0.777 \\
$\approx$v3 (ablazione config)           & 0.956 & ---   & 0.913 & ---   & 0.334 & 0.744 \\
v4 $+$ \texttt{dev:guest}                & 0.961 & ---   & 0.911 & 0.632 & 0.488 & 0.787 \\
\bottomrule
\end{tabular}
\caption{Quadro riepilogativo delle prestazioni di tutti i modelli e varianti osservati. \textbf{Solo i confronti intra-pannello sono validi} (protocolli di valutazione diversi). \textit{Lat.\ Recall} è a soglia globale 1\% FPR nel Pannello~A e a decisione instradata nei Pannelli~B e~C. La colonna \textit{agg Recall} (recall aggregata su tutte le classi) è alla soglia globale 1\% FPR nel Pannello~A e alla decisione instradata nei Pannelli~B/C. Nel Pannello~A i valori sono medie su \textbf{3 seed} \texttt{[42,7,123]} (dispersione nelle Tab.~\ref{tab:baselines} e~\ref{tab:v3v4}): la riga \textbf{TGN} è la run \emph{v3} (keying per-cookie), le baseline sono sotto la configurazione \emph{deployable} corrente (\texttt{dev:guest}); differiscono da quelli del Pannello~B per la diversa regola di decisione (soglia globale vs instradata). I valori di Pannello~C sono medie su 3 \texttt{seed} su uno stream \emph{theft-rich} (incidenti più frequenti, non appaiati evento-per-evento), perciò non comparabili in assoluto con i Pannelli~A/B. $^{\dagger}$ Il recall della GNN statica è \emph{spurio} (AUC$\approx$caso a soglia permissiva, non segnale reale). $^{\ddagger}$ XGBoost è \emph{supervisionato} (vede le etichette): upper-bound di riferimento, fuori dal paradigma non supervisionato del TGN. Lo \textit{sweep} architetturale (Tab.~\ref{tab:archsweep}) conferma inoltre che nessun aumento di capacità migliora il movimento laterale. Sorgenti: Tab.~\ref{tab:baselines} (A), Tab.~\ref{tab:v3v4} (B), Tab.~\ref{tab:theft} e~\ref{tab:guestdev} (C).}
\label{tab:summary}
\end{table}
